Este trabajo caracteriza un \textbf{filtro pasabajos RLC de segundo orden} de
frecuencia natural nominal \SI{1073}{\hertz}, y contrasta las tres fuentes con
las que se puede describir un circuito:

\begin{enumerate}[itemsep=2pt, topsep=4pt]
  \item el \textbf{modelo teórico}, evaluado a partir de la transferencia
        analítica;
  \item la \textbf{simulación} en \textsc{ngspice} del netlist, que incorpora la
        resistencia real del bobinado del inductor;
  \item la \textbf{medición} de laboratorio: barrido de respuesta en frecuencia
        y captura del transitorio en el osciloscopio.
\end{enumerate}

El objetivo no es solo verificar que las tres coinciden --- coinciden --- sino
\emph{medir en cuánto no lo hacen} y explicar por qué. La resistencia serie del
inductor y la tolerancia de los componentes bajan el factor de mérito real por
debajo del ideal, y eso se ve tanto en el pico del Bode como en el sobrepico de
la respuesta al escalón.

\begin{quote}\small
\textbf{Sobre los datos.} Las curvas rotuladas como \emph{medidas} son
sintéticas: las genera \texttt{fuentes/generar\_mediciones.py} a partir del
modelo del circuito real, con tolerancias y ruido de instrumento. Este informe
es el ejemplo de referencia de Xtal, y existe para que se pueda reproducir
entero sin tener el circuito armado en un banco.
\end{quote}
